\begin{abstract}
Battery-powered river water-quality nodes lose data, yet retransmission, buffering and reconstruction are analysed for short, independent losses. This paper measures loss in four real systems, two water-quality deployments and two LoRaWAN networks. Losses separate into short radio losses and rarer long delivery gaps, which hold most of the missing data in the water-quality deployments; a geometric gap-length model fits worst on every record. Replayed through a measurement-parameterised LoRaWAN and NB-IoT energy model, retransmission recovers under a fifth of the missing readings even under optimistic retry timing, at \arqRatioMin--\arqRatioMax{} the energy of unconfirmed reporting on LoRaWAN, whereas buffering with back-off recovers every generated reading at an energy close to that of recovering nothing. Neither simple nor deep reconstruction keeps most day-long gaps within tolerance. Bounded event-driven batching combines hourly batches, immediate alarms, a bounded-error trigger and buffering with back-off. Tested with frozen rules on held-out data from \confSites{} river sites, it matches per-reading reporting on alarm latency and real-time picture, while per-reading retransmission uses \arqOverBoundedMinConf--\arqOverBoundedMaxConf{} as much energy. Recovery of many nodes after a shared outage needs randomised probes, and large SF12 cells are limited by the gateway's acknowledgements.
\end{abstract}

\begin{IEEEkeywords}
Delivery gaps, network outages, store-and-forward, ARQ, LoRaWAN, NB-IoT, energy efficiency, water-quality monitoring, data reconstruction, event detection.
\end{IEEEkeywords}
